% Môn: Vật lý
% Lớp: 12
% Chương: Chương I. Vật lí nhiệt
% Bài: L12C1 Bài 4. Nhiệt dung riêng · Tính tốc độ tăng nhiệt độ
% ID: L12C1B3-121-TN
% ID: L12C1B3-121-TN
% Mức: NB
% ID: L12C1B3-121-TN


% Mức: VD
% Mức: VD
% ID: L12C1B3-135-TLN
% ID: L12C1B3-135-TLN
% Mức: VD
% ID: L12C1B3-135-TLN
% ID: L12C1B3-135-TLN
% Mức: VD
% ID: L12C1B4-76-TLN
%=== Câu 15 ===%
\dangbt{Tính tốc độ tăng nhiệt độ}
\begin{ex}
% ID: L12C1B4-76-TLN
% Mức: VD
Trong khoảng thời gian $2$ phút $12$ giây, nhiệt độ của một vật tăng từ $-15^\circ C$ đến $8,6^\circ C$. Nhiệt độ trung bình trong khoảng thời gian nói trên đã tăng bao nhiêu Kelvin/giây $(K/s)$ (làm tròn kết quả đến chữ số hàng phần trăm)?
	\shortans[oly]{$0,18$}%Thay đáp án chuẩn 
	\loigiai{
	Ta có: $\dfrac{\Delta T}{\Delta \tau} = \dfrac{\Delta t}{\Delta \tau} = \dfrac{8,6 - (-15)}{2.60+12} \approx 0,18\ (K/s)$.
	}
\end{ex}
